\subsection{环}

定义 1. 环是一个集合 A，带有两个合成律，分别称为加法和乘法，满足以下公理：

（AN I）在加法下，A 是一个交换群。

（AN II）乘法是结合的且具有单位元。

（AN III）乘法对加法满足分配律。

环 A 称为交换环，如果其乘法是交换的。

在下文中， \((x,y)\mapsto x+y\) 表示加法， \((x,y)\mapsto xy\) 表示乘法；0 表示加法的单位元，1 表示乘法的单位元。最后，-x 表示 x 在加法下的负元。因此环的公理由以下恒等式表达：

\[
x + (y + z) = (x + y) + z \quad (\text { associativity   of   addition })\tag{1}
\]

\[
x + y = y + x \quad (\text { commutativity   of   addition })\tag{2}
\]

\[
0 + x = x + 0 = x \quad (\text { zero })\tag{3}
\]

\[
x + (- x) = (- x) + x = 0 \quad (\text { negative })\tag{4}
\]

\[
x (y z) = (x y) z \quad (\text { associativity   of   multiplication })\tag{5}
\]

\[
x. 1 = 1. x = x \quad (\text { unit   element })\tag{6}
\]

\[
\left. \begin{array}{l} (x + y). z = x z + y z \\ x. (y + z) = x y + x z \end{array} \right\} (\text {distributivity})\tag{7, 8}
\]

最后，环 A 是交换的，如果 \(xy = yx\) 对 \(x, y\) （在 A 中）成立。

仅就加法而言，A 是一个交换群，称为 A 的加法群。对所有 \(x \in A\) ，我们定义左位似 \(\gamma_{x}\) 和右位似 \(\delta_{x}\) 为 \(\gamma_{x}(y) = xy\) , \(\delta_{x}(y) = yx\) 。由公式 (7) 和 (8)， \(\gamma_{x}\) 和 \(\delta_{x}\) 是 A 的加法群的自同态，因此把零映到零，把负元映到负元。因此

\[
x. 0 = 0. x = 0\tag{9}
\]

\[
x. (- y) = (- x). y = - x y;\tag{10}
\]

由此得出 \((-x)(-y) = -((-x).y) = -(-xy)\) ，从而

\[
(- x) (- y) = x y.\tag{11}
\]

公式 (10) 和 (11) 构成符号法则。由此得出

\[
- x = (- 1) x = x (- 1)
\]

和 \((-1)(-1) = 1\) 。

由 (11) 通过对 \(n\) 进行归纳可得

\[
(- x) ^ {n} = \left\{ \begin{array}{l l} x ^ {n} & \text { if   } n \text {   is   even } \\ - x ^ {n} & \text { if   } n \text {   is   odd. } \end{array} \right.\tag{12}
\]

当我们谈到环 A 中的可消元、可逆元、可置换元、中心元、中心化子或中心时，所有这些概念均指 A 的乘法。若 \(x, y \in A\) 且 y 是可逆的，则 A 的元素 \(xy^{-1}\) 当 A 是交换环时也记作 x/y。A 的可逆元素集在乘法下是稳定的。在乘法诱导的运算下，它构成一个群，称为 A 的乘法群，有时记为 A*。

设 \(x, y\) 属于 A 。 \(x\) 称为 \(y\) 的左（相应地，右）倍元，如果存在 \(y' \in \mathbf{A}\) 使得 \(x = y'y\) （相应地 \(x = yy'\) ）；也常说 \(y\) 是 x 的右（相应地，

左）因子。当 A 是交换的时，无需区分“左”和“右”。

按照上述术语，每个元素 \(y \in A\) 都将被视为 0 的右因子和左因子；但以语言上的滥用，术语“0 的右（相应地，左）因子”通常保留给这样的元素 y：存在 \(x \neq 0\) （在 A 中）满足关系 xy = 0（相应地，yx = 0）。换言之，右（相应地，左）零因子正是右（相应地，左）不可消去的元素。

设 \(x \in A\) 。 \(x\) 称为幂零的，如果存在整数 \(n > 0\) 使得 \(x^n = 0\) 。元素 \(1 - x\) 此时可逆，其逆等于

\[
1 + x + x ^ {2} + \dots + x ^ {n - 1}.
\]

由于 A 在加法下构成交换群，元素 \(nx\) 对于 \(n \in \mathbf{Z}\) 和 \(x \in \mathbf{A}\) 已有定义（§ 2，第 8 号）。由于 \(\gamma_x\) 和 \(\delta_x\) 是加法群 \(\mathbf{A}\) 的自同态， \(\gamma_x(ny) = n\gamma_x(y)\) 和 \(\delta_y(nx) = n\delta_y(x)\) ，因此

\[
x. (n y) = (n x). y = n. (x y).
\]

特别地， \(nx = (n.1)x\) 。

一个集合 A，带有加法和乘法，满足环的公理，但除去保证乘法单位元存在的公理，称为伪环。

